
\section{Individual and block time steps}
\label{ch:timesteps}

Stellar systems are characterized by a huge dynamical range in
radial and temporal scales.
The time scale varies e.g.\ in a star cluster from orbital periods
of binaries of some days up to the relaxation of a few hundred
million years, or even billions of years.
Even if we put for a moment the very close binaries aside,
which are treated differently (by regularization methods),
there typically is a large dynamic range in the average local
stellar density from its centre to the very outskirts, where it
dissolves into the galactic tidal field.
In a classical picture, the two closest bodies would determine
the time--step of force calculation for the whole rest of the
system.
However, for bodies in regions where the changes of the force are
relatively small, a permanent re--computing of the terms appears
time consuming.
So, in order to economize the calculation, these objects shall
be allowed to move a longer distance before a recomputation
becomes necessary.
In between there is always the possibility to acquire particle
positions and velocities via a Taylor series prediction,
as described in Chapter~\ref{ch:hermite}.
This is the idea of a vital method for assigning different
time--steps, $\Delta t = t_1 - t_0$, between the force
computations, the so--called ``individual time--step scheme''
\cite{a_63}, which was later advanced to the hierarchical
block steps.


%%%%%%%%%%%%   Figure for block time step scheme   %%%%%%%%%%%%
\begin{figure}[ht]
%
\centering
\framebox{
\unitlength1cm
\begin{picture}(10,6.2)
   % lower horizontal line (time step)
   \thinlines
      \put(0.5,1){\line(1,0){9.0}}                 % time axis
      \multiput(0.5,1.125)(0.5,0){18}{\line(0,-1){0.25}} % tickmarks
      \put(0.43,0.5){0}                             % ticknames
      \put(0.93,0.5){1}
      \put(1.43,0.5){2}
      \put(2.43,0.5){4}
      \put(4.43,0.5){8}
      \put(5.8,0.5){time steps}
      \put(8.3,0.5){16}
   % time steps of particles
   \thicklines
      \multiput(0.5,2.2)(0.5,0){12}{\vector(1,0){0.5}}
          \multiput(6.5,2.2)(1.0,0){2}{\vector(1,0){1.0}}
      \multiput(0.5,3.2)(1.0,0){4}{\vector(1,0){1.0}}
          \multiput(4.5,3.2)(2.0,0){2}{\vector(1,0){2.0}}
      \multiput(0.5,4.2)(2.0,0){2}{\vector(1,0){2.0}}
          \multiput(4.5,4.2)(1.0,0){2}{\vector(1,0){1.0}}
          \multiput(6.5,4.2)(0.5,0){4}{\vector(1,0){0.5}}
      \multiput(0.5,5.2)(4.0,0){2}{\vector(1,0){4.0}}
   % vertical dotted lines
      \put(0.5,1.8){\line(0,1){3.8}}
      \dottedline{0.2}(1.0,1)(1.0,2.1)
      \dottedline{0.2}(1.5,1)(1.5,3.1)
      \dottedline{0.2}(2.0,1)(2.0,2.1)
      \dottedline{0.2}(2.5,1)(2.5,4.1)
      \dottedline{0.2}(3.0,1)(3.0,2.1)
      \dottedline{0.2}(3.5,1)(3.5,3.1)
      \dottedline{0.2}(4.0,1)(4.0,2.1)
      \dottedline{0.2}(4.5,1)(4.5,5.1)
      \dottedline{0.2}(5.0,1)(5.0,2.1)
      \dottedline{0.2}(5.5,1)(5.5,2.1)
          \dottedline{0.2}(5.5,3.7)(5.5,4.6)
          \dottedline{0.2}(7.0,3.7)(7.0,4.6)
          \dottedline{0.2}(7.5,3.7)(7.5,4.6)
          \dottedline{0.2}(8.0,3.7)(8.0,4.6)
      \dottedline{0.2}(6.0,1)(6.0,2.1)
      \dottedline{0.2}(6.5,1)(6.5,4.1)
%      \dottedline{0.2}(7.0,1)(7.0,2.1)
      \dottedline{0.2}(7.5,1)(7.5,2.1)
%      \dottedline{0.2}(8.0,1)(8.0,2.1)
      \dottedline{0.2}(8.5,1)(8.5,5.1)
   % vertical ticknames
      \put(0.0,2.1){$i$}
      \put(0.0,3.1){$k$}
      \put(0.0,4.1){$l$}
      \put(0.0,5.1){$m$}
      \put(0.0,5.8){particles}
%      \put(0.0,5.7){blocks}
\end{picture}
}%
  \vspace{1ex}
  \caption{Block time steps exemplary for four
           particles.}
  \label{pic:its}
\end{figure}


Each particle is assigned its own $\Delta t_i$ which is
first illustrated for the case of ``block time--steps''
in Figure \ref{pic:its}.
The particle named $i$ has the smallest time step at the
beginning, so its phase space coordinates are determined at
each time step.
The time step of $k$ is twice as large as $i$'s,
and its coordinates are just extrapolated (``predicted'')
at the odd time steps, while a full force calculation is
due at the dotted times.
The step width may be altered or not after the end of the
integration cycle for the special particle, as demonstrated
for $k$ and $l$ beyond the label ``8''.
The time steps have to stay commensurable with both,
each other as well as the total time, such that a
hierarchy is guaranteed.
This is the block step scheme.

As a first estimate, the rate of change of the acceleration
seems to be a reasonable quantity for the choice of the time
step:
$\Delta t_i \propto \sqrt{ \mathbf a_i / \dot{\mathbf a}_i }$.
But it turns out that for special situations in a many-body
system, it provides some undesired numerical errors.
After some experimentation, the following formula was adopted
\cite{a_85}:
%
\begin{eqnarray}
\Delta t_i = \sqrt{
           \eta \frac{|\mathbf a_{1,i}||\mathbf a_{1,i}^{(2)}| +
                      |\dot{\mathbf a}_{1,i}|^2 }
                     {|\dot{\mathbf a}_{1,i}||\mathbf a_{1,i}^{(3)}| +
                      |\mathbf a_{1,i}^{(2)}|^2 }
                  } ,    \label{eq:timestep}
\end{eqnarray}
%
where $\eta$ is a dimensionless accuracy parameter which
controls the error.
In most applications it is taken to be $\eta \approx$ 0.01
to 0.02, see also next chapter.

For the block--time steps, the synchronization is made
by taking the next--lowest integer of $\Delta t_i$;
the time steps are quantized to powers of 2 \cite{m_91a}.
Then, there will be a group (block) of several particles
which are due to movement at each time step.
If one keeps the exact $\Delta t_i$'s evaluated from
(\ref{eq:timestep}) for each particle, the commensurability
is destroyed, and we arrive at the so--called ``individual
time steps'';
in this case, there exists one sole particle being due.
The latter concept is realized in the earlier codes
NBODY1, NBODY3, NBODY5, where a neighbour scheme is renounced.
NBODY4, NBODY6, and NBODY6++ use a block step scheme.

% In the running code, the time--steps are adjusted to their
% appropriate values fairly quick. % Makino & Aarseth (1992)
% Although successive steps normally change smoothly,
% it is prudent to restrict the growth by a stability factor
% of 1.2. % Aarseth (1985)

Subsystems like star binaries, triples or a similar subgroups
(they are termed KS pairs, chains, hierarchies) enter the
time--step scheme with their respective centre's of masses only.
Their internal motion is treated in a different way by a
regularized integration (Chapter~\ref{ch:regularization}).


